\chapter{Literature Review}
\label{ch:literature}

This chapter surveys the prior work that underpins the design
decisions taken in the remainder of the report. Section~\ref{sec:lit-physio}
reviews the physiological signals that are known to co-vary with
psychological stress and that are accessible from a consumer wrist
wearable. Section~\ref{sec:lit-datasets} compares the public datasets
available for supervised stress detection and explains why WESAD was
selected for this project. Section~\ref{sec:lit-ml} surveys the
machine-learning approaches applied to the stress-detection problem
and positions the proposed bidirectional LSTM with attention against
them. Section~\ref{sec:lit-wearables} examines commercial and research
wearables that ship stress-related features, and Section~\ref{sec:lit-cabin}
examines adaptive vehicle cabin systems. Finally, Section~\ref{sec:lit-gaps}
consolidates the four gaps identified in Chapter~\ref{ch:introduction}
and anchors each of them to the reviewed literature.

\section{Physiological Correlates of Stress}
\label{sec:lit-physio}

Psychological stress is mediated by the autonomic nervous system
and produces measurable changes in peripheral physiology.
\citet{sharma2012driver} survey more than two decades of sensor
technologies and computational techniques for stress recognition.
Among the many physiological signals studied, four are both strongly
predictive of stress and accessible from a wrist-worn consumer device.

\textbf{Heart-rate variability (HRV)} derived from the photoplethysmography
(PPG) signal is the most widely studied. When stress hits, the
parasympathetic branch of the autonomic nervous system withdraws its
tonic inhibition of the sinoatrial node and heart rate variability
collapses. Time-domain indices like RMSSD and SDNN both drop during
acute stressors, and frequency-domain indices shift toward the
low-frequency band (0.04--0.15~Hz) at the expense of the high-frequency
band (0.15--0.4~Hz). These changes can be tracked reliably from
inter-beat intervals (IBIs) extracted from wrist PPG.

\textbf{Electrodermal activity (EDA)}, also called galvanic skin response,
reflects eccrine sweat-gland activity controlled exclusively by the
sympathetic nervous system. Under stress, the tonic component of EDA
rises and the density and amplitude of phasic peaks per minute
increases. EDA is only weakly coupled to physical workload, which makes
it one of the cleaner peripheral indicators of psychological stress;
it is a key signal in both \citet{schmidt2018wesad} and
\citet{can2019stress}. Not all consumer wearables ship a dedicated
EDA sensor, but the WESAD protocol was designed around wrist EDA and
provides a labeled corpus for supervised learning.

\textbf{Skin temperature} responds to sympathetic activation through
peripheral vasoconstriction, which diverts blood flow away from the
distal extremities. Within a few minutes of a stressor, wrist skin
temperature typically drops 0.1--0.5~$^\circ$C below baseline. The
absolute change is small compared to EDA, but the signal is cheap
to measure, drifts slowly, and is compatible with the 4~Hz sampling
rate that consumer wearables expose. It's present on virtually every
modern smartwatch.

\textbf{Three-axis acceleration} doesn't encode stress directly --- it
encodes movement. But movement is a critical confound because both PPG
and EDA are contaminated by motion artifacts. Including the accelerometer
magnitude as a feature lets the model learn to down-weight signals that
coincide with high-magnitude motion, which is why it's included as the
fifth input feature.

Together, these four channels cover both branches of the autonomic
nervous system (sympathetic via EDA and temperature,
parasympathetic-sensitive via HRV), one direct cardiovascular channel
(PPG), and one confound-control channel (acceleration) --- all from a
modern wrist wearable. This is the five-feature input vector used
throughout the rest of this report.

\section{Existing Stress Detection Datasets}
\label{sec:lit-datasets}

Supervised stress detection requires labeled multimodal physiological
data collected under controlled conditions. Several public datasets
have been released in the past decade; the three most relevant to
this project are summarized in Table~\ref{tab:datasets}.

\begin{table}[H]
\centering
\caption{Comparison of public multimodal stress-detection datasets
relevant to this work.}
\label{tab:datasets}
\small
\begin{tabular}{lcp{2.5cm}p{4.0cm}c}
\toprule
\textbf{Dataset} & \textbf{Subjects} & \textbf{Sensors (wrist)} &
\textbf{Conditions} & \textbf{Classes} \\
\midrule
WESAD \citep{schmidt2018wesad} & 15 & BVP/EDA/TEMP/ACC & Baseline, Stress,
Amusement, Meditation & 3 \\
DEAP & 32 & EEG+EDA+Resp & Video-induced emotion & Valence/Arousal \\
SWELL-KW \citep{koldijk2014swell} & 25 & ECG+EDA & Office tasks,
time pressure & 3 \\
\bottomrule
\end{tabular}
\end{table}

\textbf{WESAD} \citep{schmidt2018wesad} is the most widely used
benchmark in the wearable stress-detection community. It contains
synchronized recordings from fifteen subjects (S2--S11 and S13--S17;
S12 is missing from the public release) wearing both a RespiBAN
chest device and an Empatica E4 wrist device. Each subject underwent
a fixed protocol consisting of baseline, amusement (video clips),
stress (a Trier Social Stress Test variant featuring public speaking
and mental arithmetic), and meditation. Sampling rates for the wrist
sensors are 64~Hz for BVP, 4~Hz for EDA and skin temperature,
32~Hz for three-axis acceleration, and 700~Hz for the event-label
channel. This combination of sensor availability and protocol
structure is ideally matched to our four-class formulation, and it
is the only one of the three datasets that uses \emph{exactly} the
same wrist sensors that we target on the deployment device.

\textbf{DEAP} is a large emotion-recognition dataset, but it is
primarily EEG-based and its labels are continuous valence/arousal
ratings rather than discrete stress categories. EEG is not a wrist
modality, so DEAP is of limited relevance to on-device deployment.

\textbf{SWELL-KW} \citep{koldijk2014swell} is an office-workload
dataset focused on knowledge-work stressors (interruptions, time
pressure). It uses ECG and EDA collected via laboratory-grade
equipment. While its label structure is closer to ours than DEAP's,
it does not include BVP or skin temperature, and its stressor
profile (cognitive, sustained) is different from the acute,
arousal-dominated stressors that apply to driving.

For these reasons WESAD is the primary dataset used throughout the
remainder of this report. SWELL-KW is retained only as a
cross-reference in Chapter~\ref{ch:results}.

\section{Machine Learning for Stress Classification}
\label{sec:lit-ml}

A wide range of classifiers has been applied to the stress-detection
problem, spanning classical feature-engineering pipelines and modern
deep-learning architectures. A representative selection, together
with the approach taken in this work, is summarized in
Table~\ref{tab:ml-comparison}.

\begin{table}[H]
\centering
\caption{Representative machine-learning approaches to wearable stress
detection.}
\label{tab:ml-comparison}
\small
\begin{tabular}{p{3.3cm}p{2.3cm}p{1.5cm}p{4.5cm}}
\toprule
\textbf{Approach} & \textbf{Dataset} & \textbf{Acc.} & \textbf{Limitation} \\
\midrule
Handcrafted features + SVM \citep{schmidt2018wesad} & WESAD & $\approx$ 80\% &
Sensitive to feature selection, no temporal modelling. \\
Random forest on TSFEL features \citep{barandas2020tsfel} & WESAD &
$\approx$ 85\% & Large feature vector, not deployable on wearable. \\
CNN on raw signals \citep{can2019stress} & Multiple & $\approx$ 87\% &
Poor long-range temporal context. \\
LSTM on windowed features & WESAD & $\approx$ 88\% & Unidirectional;
misses future context within window. \\
\textbf{BiLSTM + attention (this work)} & \textbf{WESAD} & \textbf{97.8\%}
& \textbf{Wrist-only, 72~KB TFLite, real-time on Wear~OS.} \\
\bottomrule
\end{tabular}
\end{table}

The earliest approaches applied support vector machines (SVMs) to
handcrafted statistical features extracted from short windows of
physiological signals. These methods, exemplified by the baseline
classifier reported in the WESAD paper itself \citep{schmidt2018wesad},
are simple to train but brittle: they rely on dataset-specific
feature engineering and cannot capture the temporal structure that
a bidirectional recurrent network exploits naturally. Libraries such
as TSFEL \citep{barandas2020tsfel} have automated the feature-extraction
stage, but the resulting pipelines still depend on large feature
vectors that are expensive to compute on-device.

Convolutional neural networks (CNNs) were the next natural step.
By treating the multi-channel physiological signal as a one-dimensional
``image'', CNNs can learn local patterns (e.g.\ PPG waveform shape,
EDA peak morphology) without hand-designed features. However, CNNs
alone capture only local temporal context: they cannot model the
slow, asymmetric dynamics of stress onset and recovery that take
place over tens of seconds.

Recurrent networks, and in particular \textbf{long short-term memory}
(LSTM) networks, overcome this limitation by maintaining an explicit
hidden state across the full window. \citet{graves2005bilstm} showed
that \emph{bidirectional} LSTMs -- which run one recurrent pass forward
in time and another backward -- further improve classification accuracy
on sequence-labelling problems because each step of the hidden state
has access to both past and future context within the window. This is
particularly relevant for stress classification because physiological
responses are often delayed relative to the stressor event, and
looking ``backward'' within a thirty-second window allows the network
to associate a late EDA peak with its earlier trigger.

The remaining difficulty is that a bidirectional LSTM produces a
sequence of hidden states, one per timestep, and the downstream
classifier must collapse this sequence into a single vector. The
simplest strategies -- mean pooling, last-state pooling -- treat every
timestep as equally informative, which is demonstrably wrong for
stress detection: the informative content is concentrated in the
few timesteps where EDA peaks or where the inter-beat interval
shortens. \citet{vaswani2017attention} introduced the scaled
dot-product attention that revolutionised sequence modelling; in
this work we use a much simpler additive attention mechanism (a
single dense layer followed by a softmax across timesteps), which
is sufficient to weight informative timesteps above uninformative
ones while costing almost nothing in parameters.

Our final architecture, presented in Chapter~\ref{ch:methodology},
therefore combines a bidirectional LSTM backbone with an additive
attention head. The appended row of Table~\ref{tab:ml-comparison}
summarizes its performance: 97.8\% accuracy on held-out WESAD data
at a deployment footprint of 72~KB, which is roughly four times
smaller than the unquantized baseline and hundreds of times smaller
than CNN-based alternatives in the literature.

\section{Wearable Stress Monitoring Systems}
\label{sec:lit-wearables}

Commercial wearables have increasingly advertised ``stress tracking''
as a native feature. The Fitbit Sense, the Apple Watch (via the
Mindfulness app and HRV notifications), the Garmin Venu, and the
Samsung Galaxy Watch all report a stress score to the user. However,
these scores are generally computed from HRV alone, over minute-scale
windows, and they are neither open nor validated against a public
ground-truth dataset. They are also not exposed to third-party
applications in a form that can be used for automotive actuation.

On the research side, \citet{can2019stress} survey more than forty
published smart-phone and wearable stress-detection systems. The majority target office or
ambulatory settings rather than driving, and very few report a
deployed end-to-end pipeline from wearable sensing to a downstream
actuator. The Empatica E4 and the Shimmer3 are frequently used as
research-grade wearables, but their acquisition cost and form factor
place them outside the reach of mainstream deployment.

The gap that our system addresses is therefore twofold: (i) it uses
a consumer-grade wrist configuration compatible with Wear OS devices
in the hundreds-of-dollars price range, and (ii) it exposes its
inference output through a documented BLE characteristic, allowing
any BLE-capable vehicle module to consume it.

\section{Adaptive Vehicle Environment Systems}
\label{sec:lit-cabin}

Modern vehicles already provide some degree of cabin-environment
personalisation. Ambient LED lighting, multi-zone climate control,
configurable audio equalisation, and seat-position memory are all
standard features of mid-range and premium vehicles. Several
manufacturers have also introduced static ``driving modes'' (``Comfort'',
``Sport'', ``Eco'') that adjust chassis and powertrain parameters, and
a smaller number of manufacturers offer ``mood'' presets that group
together LED color, audio, and climate settings.

What these systems have in common is that \emph{they are all configured
manually by the driver}. None of them close the loop with the driver's
real-time physiological state; none of them adapt automatically as the
driver's stress level rises or falls; and none of them are exposed to
third-party physiological inputs. The academic literature has proposed
closed-loop cabin-control systems in the context of automated driving,
but these proposals are typically simulator-based, rely on camera-based
emotion recognition rather than physiological sensing, and stop short
of a deployable prototype.

Our system sits in a gap that no existing product or published prototype
fills: a physiologically driven, graduated, real-time cabin-environment
controller running on commodity hardware, decoupled from the vehicle's
primary electronic control units.

\section{Gap Analysis}
\label{sec:lit-gaps}

Revisiting the four gaps from Section~\ref{sec:problem} in light of
the literature reviewed above:

The \textbf{first gap --- no real-time multimodal wearable stress
classifier for drivers} --- remains open. \citet{schmidt2018wesad}
provides the dataset and \citet{can2019stress} surveys classification
approaches, but no published work combines a wrist-only feature set,
a four-class ordinal output, and an on-device deployment under 80~KB.

The \textbf{second gap --- no adaptive cabin response system} ---
is equally unaddressed. Section~\ref{sec:lit-cabin} showed that all
existing cabin personalization systems are configured manually or
driven by camera-based emotion recognition, not by physiological
sensing.

\textbf{Existing solutions are not wearable-constrained} (gap three).
Pipelines based on TSFEL \citep{barandas2020tsfel} or CNNs
\citep{can2019stress} assume memory and compute budgets orders of
magnitude beyond what a Wear~OS device offers.

\textbf{No integrated end-to-end open system exists} (gap four).
Neither the commercial wearables reviewed in Section~\ref{sec:lit-wearables}
nor the vehicle systems reviewed in Section~\ref{sec:lit-cabin}
expose a documented interface across the wearable/vehicle boundary.

The remaining chapters address each of these gaps directly.
